% --- Archon physics formalization source begin ---
% archon:physics
% archon:covers PhyXMiniProblems/problem_phyx_mini_0760.lean
% archon:source-report reports/phyx_mini/problem_phyx_mini_0760.source.json

\chapter{Physics problem phyx\_mini\_0760}
\label{ch:PhyXMiniProblems_problem_phyx_mini_0760}

\paragraph{Problem source.}
\#\# Physical scenario

In earlier days, horses pulled barges down canals in the manner shown in figure. Suppose the horse pulls on the rope with a force of \$7900\textbackslash{} N\$ at an angle of \$\textbackslash{}theta = 18\textasciicircum{}\textbackslash{}circ\$ to the direction of motion of the barge, which is headed straight along the positive direction of an \$x\$ axis. The mass of the barge is \$9500\textbackslash{} kg\$, and the magnitude of its acceleration is \$0.12\textbackslash{} m/s\textasciicircum{}2\$.

\#\# Question

What is direction (relative to positive \$x\$) of the force on the barge from the water?

\#\# Answer choices

- A: A: \$188\textasciicircum{}\textbackslash{}circ\$

- B: B: \$194\textasciicircum{}\textbackslash{}circ\$

- C: C: \$201\textasciicircum{}\textbackslash{}circ\$

- D: D: \$209\textasciicircum{}\textbackslash{}circ\$

\#\# Figure caption (auxiliary; use the image as primary evidence)

The image depicts a scene involving a boat or floating object in water. Here's a detailed description:

1. **Water and Object**: A large rectangular object is floating in what appears to be a body of water. The water is illustrated with light blue color, and there are wavy lines indicating the flow and movement of water around the object.

2. **Rope and Anchor**: A rope is attached to the front of the floating object, extending out of the water at an angle. At the other end of the rope, there's an anchor resting on the seabed.

3. **Angle Indicator**: There is an angle marked as θ between the rope and the horizontal line above the water, indicating the direction of the rope relative to the object and water surface.

4. **Flow Lines**: The dashed lines around the object suggest water flow patterns, showing how the water moves around the object as it floats.

Overall, the image seems to depict a basic physics or engineering scenario related to anchoring or fluid dynamics, with the water's motion and the angle of anchoring depicted clearly.

\#\# Dataset metadata

Dynamics; Physical Model Grounding Reasoning, Multi-Formula Reasoning

\paragraph{Recorded answer/context.}
C: C: \$201\textasciicircum{}\textbackslash{}circ\$

\paragraph{Figure/image path.}
/root/proposal\_for\_physic/science-mango/phyx\_mini\_run/phyx\_data/test\_image/760.png

\paragraph{Formalization target.}
create a compiling Lean file with sorry bodies at `PhyXMiniProblems/problem\_phyx\_mini\_0760.lean`. The Lean declarations must preserve the physical quantities, dimensions or dimensional roles, figure labels, governing-law hypotheses, and final relation expressed by this problem.
Use Mathlib/Physlib names found through LeanExplore where available. If a physics API is missing, introduce faithful local abstractions rather than scalar placeholder aliases.

\begin{theorem}[Physics formalization target]
\label{thm:physics:phyx_mini_0760:target}
\lean{PhyXMiniProblems.ProblemPhyXMini0760.problem_phyx_mini_0760}
\uses{def:physics:phyx-mini-0760:phyxminiproblems-problemphyxmini0760-degreestoradians, def:physics:phyx-mini-0760:phyxminiproblems-problemphyxmini0760-anglefromdegrees, def:physics:phyx-mini-0760:phyxminiproblems-problemphyxmini0760-canalbargesetup, def:physics:phyx-mini-0760:phyxminiproblems-problemphyxmini0760-matchesproblemreadouts, def:physics:phyx-mini-0760:phyxminiproblems-problemphyxmini0760-matchesprimarycanalfigure, def:physics:phyx-mini-0760:phyxminiproblems-problemphyxmini0760-hasphysicalbargeparameters, def:physics:phyx-mini-0760:phyxminiproblems-problemphyxmini0760-satisfiesstraightlinekinematics, def:physics:phyx-mini-0760:phyxminiproblems-problemphyxmini0760-satisfiesplanarforcedirectionlaw, def:physics:phyx-mini-0760:phyxminiproblems-problemphyxmini0760-satisfiesbargenewtonssecondlaw, def:physics:phyx-mini-0760:phyxminiproblems-problemphyxmini0760-isuniqueclosestdisplayeddirection}
The assigned autoformalize agent should translate this physics problem into Lean declarations in the covered file, with theorem and lemma proof bodies written as `by sorry`.
\end{theorem}
\begin{proof}
This is an autoformalization task, not a proof task. Produce statements that can later be proved by the physics prover without weakening the physical model.
\end{proof}

% --- Archon declaration topology begin ---
\section{Declaration topology}
\begin{definition}[Planar Vector]
  \label{def:physics:phyx-mini-0760:phyxminiproblems-problemphyxmini0760-planarvector}
  \lean{PhyXMiniProblems.ProblemPhyXMini0760.PlanarVector}
  Two-dimensional Euclidean vectors in the horizontal canal plane
\end{definition}
\begin{proof}
  This is the declared model definition of Planar Vector.
\end{proof}
\begin{definition}[acceleration Dimension]
  \label{def:physics:phyx-mini-0760:phyxminiproblems-problemphyxmini0760-accelerationdimension}
  \lean{PhyXMiniProblems.ProblemPhyXMini0760.accelerationDimension}
  The physical dimension L T⁻² of acceleration
\end{definition}
\begin{proof}
  This is the declared model definition of acceleration Dimension.
\end{proof}
\begin{definition}[force Dimension]
  \label{def:physics:phyx-mini-0760:phyxminiproblems-problemphyxmini0760-forcedimension}
  \lean{PhyXMiniProblems.ProblemPhyXMini0760.forceDimension}
  The physical dimension M L T⁻² of force
\end{definition}
\begin{proof}
  This is the declared model definition of force Dimension.
\end{proof}
\begin{definition}[Mass Quantity]
  \label{def:physics:phyx-mini-0760:phyxminiproblems-problemphyxmini0760-massquantity}
  \lean{PhyXMiniProblems.ProblemPhyXMini0760.MassQuantity}
  A nonnegative physical mass, independent of the unit used to read it
\end{definition}
\begin{proof}
  This is the declared model definition of Mass Quantity.
\end{proof}
\begin{definition}[Planar Acceleration Quantity]
  \label{def:physics:phyx-mini-0760:phyxminiproblems-problemphyxmini0760-planaraccelerationquantity}
  \lean{PhyXMiniProblems.ProblemPhyXMini0760.PlanarAccelerationQuantity}
  \uses{def:physics:phyx-mini-0760:phyxminiproblems-problemphyxmini0760-planarvector, def:physics:phyx-mini-0760:phyxminiproblems-problemphyxmini0760-accelerationdimension}
  A signed planar physical acceleration vector
\end{definition}
\begin{proof}
  This is the declared model definition of Planar Acceleration Quantity.
\end{proof}
\begin{definition}[Planar Force Quantity]
  \label{def:physics:phyx-mini-0760:phyxminiproblems-problemphyxmini0760-planarforcequantity}
  \lean{PhyXMiniProblems.ProblemPhyXMini0760.PlanarForceQuantity}
  \uses{def:physics:phyx-mini-0760:phyxminiproblems-problemphyxmini0760-planarvector, def:physics:phyx-mini-0760:phyxminiproblems-problemphyxmini0760-forcedimension}
  A signed planar physical force vector
\end{definition}
\begin{proof}
  This is the declared model definition of Planar Force Quantity.
\end{proof}
\begin{definition}[x Axis]
  \label{def:physics:phyx-mini-0760:phyxminiproblems-problemphyxmini0760-xaxis}
  \lean{PhyXMiniProblems.ProblemPhyXMini0760.xAxis}
  Coordinate 0, the positive direction of barge motion
\end{definition}
\begin{proof}
  This is the declared model definition of x Axis.
\end{proof}
\begin{definition}[y Axis]
  \label{def:physics:phyx-mini-0760:phyxminiproblems-problemphyxmini0760-yaxis}
  \lean{PhyXMiniProblems.ProblemPhyXMini0760.yAxis}
  Coordinate 1, the side of the canal on which the horse is shown
\end{definition}
\begin{proof}
  This is the declared model definition of y Axis.
\end{proof}
\begin{definition}[x Hat]
  \label{def:physics:phyx-mini-0760:phyxminiproblems-problemphyxmini0760-xhat}
  \lean{PhyXMiniProblems.ProblemPhyXMini0760.xHat}
  \uses{def:physics:phyx-mini-0760:phyxminiproblems-problemphyxmini0760-planarvector, def:physics:phyx-mini-0760:phyxminiproblems-problemphyxmini0760-xaxis}
  The dimensionless unit vector along positive x
\end{definition}
\begin{proof}
  This is the declared model definition of x Hat.
\end{proof}
\begin{definition}[y Hat]
  \label{def:physics:phyx-mini-0760:phyxminiproblems-problemphyxmini0760-yhat}
  \lean{PhyXMiniProblems.ProblemPhyXMini0760.yHat}
  \uses{def:physics:phyx-mini-0760:phyxminiproblems-problemphyxmini0760-planarvector, def:physics:phyx-mini-0760:phyxminiproblems-problemphyxmini0760-yaxis}
  The dimensionless unit vector along positive y
\end{definition}
\begin{proof}
  This is the declared model definition of y Hat.
\end{proof}
\begin{definition}[nonnegative Readout]
  \label{def:physics:phyx-mini-0760:phyxminiproblems-problemphyxmini0760-nonnegativereadout}
  \lean{PhyXMiniProblems.ProblemPhyXMini0760.nonnegativeReadout}
  Read a nonnegative physical scalar in a coherent unit system
\end{definition}
\begin{proof}
  This is the declared model definition of nonnegative Readout.
\end{proof}
\begin{definition}[planar Vector Readout]
  \label{def:physics:phyx-mini-0760:phyxminiproblems-problemphyxmini0760-planarvectorreadout}
  \lean{PhyXMiniProblems.ProblemPhyXMini0760.planarVectorReadout}
  \uses{def:physics:phyx-mini-0760:phyxminiproblems-problemphyxmini0760-planarvector}
  Read a planar physical vector in a coherent unit system
\end{definition}
\begin{proof}
  This is the declared model definition of planar Vector Readout.
\end{proof}
\begin{definition}[mass In Kilograms]
  \label{def:physics:phyx-mini-0760:phyxminiproblems-problemphyxmini0760-massinkilograms}
  \lean{PhyXMiniProblems.ProblemPhyXMini0760.massInKilograms}
  \uses{def:physics:phyx-mini-0760:phyxminiproblems-problemphyxmini0760-massquantity, def:physics:phyx-mini-0760:phyxminiproblems-problemphyxmini0760-nonnegativereadout}
  Kilogram readout of the barge's physical mass
\end{definition}
\begin{proof}
  This is the declared model definition of mass In Kilograms.
\end{proof}
\begin{definition}[acceleration In Meters Per Second Squared]
  \label{def:physics:phyx-mini-0760:phyxminiproblems-problemphyxmini0760-accelerationinmeterspersecondsquared}
  \lean{PhyXMiniProblems.ProblemPhyXMini0760.accelerationInMetersPerSecondSquared}
  \uses{def:physics:phyx-mini-0760:phyxminiproblems-problemphyxmini0760-planarvector, def:physics:phyx-mini-0760:phyxminiproblems-problemphyxmini0760-planaraccelerationquantity, def:physics:phyx-mini-0760:phyxminiproblems-problemphyxmini0760-planarvectorreadout}
  Cartesian SI readout of acceleration, in metres per second squared
\end{definition}
\begin{proof}
  This is the declared model definition of acceleration In Meters Per Second Squared.
\end{proof}
\begin{definition}[force Vector In Newtons]
  \label{def:physics:phyx-mini-0760:phyxminiproblems-problemphyxmini0760-forcevectorinnewtons}
  \lean{PhyXMiniProblems.ProblemPhyXMini0760.forceVectorInNewtons}
  \uses{def:physics:phyx-mini-0760:phyxminiproblems-problemphyxmini0760-planarvector, def:physics:phyx-mini-0760:phyxminiproblems-problemphyxmini0760-planarforcequantity, def:physics:phyx-mini-0760:phyxminiproblems-problemphyxmini0760-planarvectorreadout}
  Cartesian SI readout of force, in newtons
\end{definition}
\begin{proof}
  This is the declared model definition of force Vector In Newtons.
\end{proof}
\begin{definition}[acceleration Magnitude In Meters Per Second Squared]
  \label{def:physics:phyx-mini-0760:phyxminiproblems-problemphyxmini0760-accelerationmagnitudeinmeterspersecondsquared}
  \lean{PhyXMiniProblems.ProblemPhyXMini0760.accelerationMagnitudeInMetersPerSecondSquared}
  \uses{def:physics:phyx-mini-0760:phyxminiproblems-problemphyxmini0760-planaraccelerationquantity, def:physics:phyx-mini-0760:phyxminiproblems-problemphyxmini0760-accelerationinmeterspersecondsquared}
  Magnitude of a planar acceleration readout, in metres per second squared
\end{definition}
\begin{proof}
  This is the declared model definition of acceleration Magnitude In Meters Per Second Squared.
\end{proof}
\begin{definition}[force Magnitude In Newtons]
  \label{def:physics:phyx-mini-0760:phyxminiproblems-problemphyxmini0760-forcemagnitudeinnewtons}
  \lean{PhyXMiniProblems.ProblemPhyXMini0760.forceMagnitudeInNewtons}
  \uses{def:physics:phyx-mini-0760:phyxminiproblems-problemphyxmini0760-planarforcequantity, def:physics:phyx-mini-0760:phyxminiproblems-problemphyxmini0760-forcevectorinnewtons}
  Magnitude of a planar force readout, in newtons
\end{definition}
\begin{proof}
  This is the declared model definition of force Magnitude In Newtons.
\end{proof}
\begin{definition}[degrees To Radians]
  \label{def:physics:phyx-mini-0760:phyxminiproblems-problemphyxmini0760-degreestoradians}
  \lean{PhyXMiniProblems.ProblemPhyXMini0760.degreesToRadians}
  Convert a degree readout to radians
\end{definition}
\begin{proof}
  This is the declared model definition of degrees To Radians.
\end{proof}
\begin{definition}[angle From Degrees]
  \label{def:physics:phyx-mini-0760:phyxminiproblems-problemphyxmini0760-anglefromdegrees}
  \lean{PhyXMiniProblems.ProblemPhyXMini0760.angleFromDegrees}
  \uses{def:physics:phyx-mini-0760:phyxminiproblems-problemphyxmini0760-degreestoradians}
  Interpret a degree readout as an angle modulo a full turn
\end{definition}
\begin{proof}
  This is the declared model definition of angle From Degrees.
\end{proof}
\begin{definition}[direction Unit Vector]
  \label{def:physics:phyx-mini-0760:phyxminiproblems-problemphyxmini0760-directionunitvector}
  \lean{PhyXMiniProblems.ProblemPhyXMini0760.directionUnitVector}
  \uses{def:physics:phyx-mini-0760:phyxminiproblems-problemphyxmini0760-planarvector, def:physics:phyx-mini-0760:phyxminiproblems-problemphyxmini0760-xhat, def:physics:phyx-mini-0760:phyxminiproblems-problemphyxmini0760-yhat}
  The Cartesian unit vector corresponding to a direction from positive x
\end{definition}
\begin{proof}
  This is the declared model definition of direction Unit Vector.
\end{proof}
\begin{definition}[Force Source]
  \label{def:physics:phyx-mini-0760:phyxminiproblems-problemphyxmini0760-forcesource}
  \lean{PhyXMiniProblems.ProblemPhyXMini0760.ForceSource}
  The two horizontal agents exerting forces on the barge
\end{definition}
\begin{proof}
  This is the declared model definition of Force Source.
\end{proof}
\begin{definition}[Figure Object]
  \label{def:physics:phyx-mini-0760:phyxminiproblems-problemphyxmini0760-figureobject}
  \lean{PhyXMiniProblems.ProblemPhyXMini0760.FigureObject}
  Physical objects visible in the supplied primary bitmap
\end{definition}
\begin{proof}
  This is the declared model definition of Figure Object.
\end{proof}
\begin{definition}[Figure Annotation]
  \label{def:physics:phyx-mini-0760:phyxminiproblems-problemphyxmini0760-figureannotation}
  \lean{PhyXMiniProblems.ProblemPhyXMini0760.FigureAnnotation}
  Literal geometric annotations visible in the supplied primary bitmap
\end{definition}
\begin{proof}
  This is the declared model definition of Figure Annotation.
\end{proof}
\begin{definition}[Rope Endpoint]
  \label{def:physics:phyx-mini-0760:phyxminiproblems-problemphyxmini0760-ropeendpoint}
  \lean{PhyXMiniProblems.ProblemPhyXMini0760.RopeEndpoint}
  The two visibly distinguished endpoints of the tow rope
\end{definition}
\begin{proof}
  This is the declared model definition of Rope Endpoint.
\end{proof}
\begin{definition}[Supplied Canal Figure]
  \label{def:physics:phyx-mini-0760:phyxminiproblems-problemphyxmini0760-suppliedcanalfigure}
  \lean{PhyXMiniProblems.ProblemPhyXMini0760.SuppliedCanalFigure}
  \uses{def:physics:phyx-mini-0760:phyxminiproblems-problemphyxmini0760-figureobject, def:physics:phyx-mini-0760:phyxminiproblems-problemphyxmini0760-figureannotation, def:physics:phyx-mini-0760:phyxminiproblems-problemphyxmini0760-ropeendpoint}
  Qualitative evidence transcribed from image 760.png. These fields record only visible objects and geometry; no water-force direction or answer value is stored in the figure
\end{definition}
\begin{proof}
  This is the declared model definition of Supplied Canal Figure.
\end{proof}
\begin{definition}[Canal Barge Setup]
  \label{def:physics:phyx-mini-0760:phyxminiproblems-problemphyxmini0760-canalbargesetup}
  \lean{PhyXMiniProblems.ProblemPhyXMini0760.CanalBargeSetup}
  \uses{def:physics:phyx-mini-0760:phyxminiproblems-problemphyxmini0760-massquantity, def:physics:phyx-mini-0760:phyxminiproblems-problemphyxmini0760-planaraccelerationquantity, def:physics:phyx-mini-0760:phyxminiproblems-problemphyxmini0760-planarforcequantity, def:physics:phyx-mini-0760:phyxminiproblems-problemphyxmini0760-forcesource, def:physics:phyx-mini-0760:phyxminiproblems-problemphyxmini0760-suppliedcanalfigure}
  Independent quantities in the barge experiment. In particular, the water force and its direction are not defined from a displayed answer; they are constrained only by the general laws below
\end{definition}
\begin{proof}
  This is the declared model definition of Canal Barge Setup.
\end{proof}
\begin{definition}[Matches Problem Readouts]
  \label{def:physics:phyx-mini-0760:phyxminiproblems-problemphyxmini0760-matchesproblemreadouts}
  \lean{PhyXMiniProblems.ProblemPhyXMini0760.MatchesProblemReadouts}
  \uses{def:physics:phyx-mini-0760:phyxminiproblems-problemphyxmini0760-massinkilograms, def:physics:phyx-mini-0760:phyxminiproblems-problemphyxmini0760-accelerationmagnitudeinmeterspersecondsquared, def:physics:phyx-mini-0760:phyxminiproblems-problemphyxmini0760-forcemagnitudeinnewtons, def:physics:phyx-mini-0760:phyxminiproblems-problemphyxmini0760-anglefromdegrees, def:physics:phyx-mini-0760:phyxminiproblems-problemphyxmini0760-canalbargesetup}
  Numerical and directional data explicitly supplied by the problem
\end{definition}
\begin{proof}
  This is the declared model definition of Matches Problem Readouts.
\end{proof}
\begin{definition}[Matches Primary Canal Figure]
  \label{def:physics:phyx-mini-0760:phyxminiproblems-problemphyxmini0760-matchesprimarycanalfigure}
  \lean{PhyXMiniProblems.ProblemPhyXMini0760.MatchesPrimaryCanalFigure}
  \uses{def:physics:phyx-mini-0760:phyxminiproblems-problemphyxmini0760-canalbargesetup}
  Primary-image evidence. The auxiliary caption's alleged anchor is omitted: the right endpoint is visibly held by the horse, and the rope rises above the right-pointing horizontal reference ray
\end{definition}
\begin{proof}
  This is the declared model definition of Matches Primary Canal Figure.
\end{proof}
\begin{definition}[Has Physical Barge Parameters]
  \label{def:physics:phyx-mini-0760:phyxminiproblems-problemphyxmini0760-hasphysicalbargeparameters}
  \lean{PhyXMiniProblems.ProblemPhyXMini0760.HasPhysicalBargeParameters}
  \uses{def:physics:phyx-mini-0760:phyxminiproblems-problemphyxmini0760-massinkilograms, def:physics:phyx-mini-0760:phyxminiproblems-problemphyxmini0760-accelerationmagnitudeinmeterspersecondsquared, def:physics:phyx-mini-0760:phyxminiproblems-problemphyxmini0760-forcemagnitudeinnewtons, def:physics:phyx-mini-0760:phyxminiproblems-problemphyxmini0760-canalbargesetup}
  Positivity conditions selecting the physical, nondegenerate setup
\end{definition}
\begin{proof}
  This is the declared model definition of Has Physical Barge Parameters.
\end{proof}
\begin{definition}[Satisfies Straight Line Kinematics]
  \label{def:physics:phyx-mini-0760:phyxminiproblems-problemphyxmini0760-satisfiesstraightlinekinematics}
  \lean{PhyXMiniProblems.ProblemPhyXMini0760.SatisfiesStraightLineKinematics}
  \uses{def:physics:phyx-mini-0760:phyxminiproblems-problemphyxmini0760-accelerationinmeterspersecondsquared, def:physics:phyx-mini-0760:phyxminiproblems-problemphyxmini0760-accelerationmagnitudeinmeterspersecondsquared, def:physics:phyx-mini-0760:phyxminiproblems-problemphyxmini0760-directionunitvector, def:physics:phyx-mini-0760:phyxminiproblems-problemphyxmini0760-canalbargesetup}
  Because the barge accelerates straight along its positive direction of motion, the acceleration vector is its stated magnitude times the corresponding unit direction. This is a kinematic modeling relation, not a water-force answer
\end{definition}
\begin{proof}
  This is the declared model definition of Satisfies Straight Line Kinematics.
\end{proof}
\begin{definition}[Satisfies Planar Force Direction Law]
  \label{def:physics:phyx-mini-0760:phyxminiproblems-problemphyxmini0760-satisfiesplanarforcedirectionlaw}
  \lean{PhyXMiniProblems.ProblemPhyXMini0760.SatisfiesPlanarForceDirectionLaw}
  \uses{def:physics:phyx-mini-0760:phyxminiproblems-problemphyxmini0760-planarvectorreadout, def:physics:phyx-mini-0760:phyxminiproblems-problemphyxmini0760-directionunitvector, def:physics:phyx-mini-0760:phyxminiproblems-problemphyxmini0760-forcesource, def:physics:phyx-mini-0760:phyxminiproblems-problemphyxmini0760-canalbargesetup}
  Each applied force decomposes into its magnitude and direction relative to positive x. The law is stated in every coherent unit system and does not assign either force a special numerical direction
\end{definition}
\begin{proof}
  This is the declared model definition of Satisfies Planar Force Direction Law.
\end{proof}
\begin{definition}[Satisfies Barge Newtons Second Law]
  \label{def:physics:phyx-mini-0760:phyxminiproblems-problemphyxmini0760-satisfiesbargenewtonssecondlaw}
  \lean{PhyXMiniProblems.ProblemPhyXMini0760.SatisfiesBargeNewtonsSecondLaw}
  \uses{def:physics:phyx-mini-0760:phyxminiproblems-problemphyxmini0760-nonnegativereadout, def:physics:phyx-mini-0760:phyxminiproblems-problemphyxmini0760-planarvectorreadout, def:physics:phyx-mini-0760:phyxminiproblems-problemphyxmini0760-canalbargesetup}
  Newton's second law in the horizontal canal plane. The water force here is the total horizontal force exerted by the surrounding water; vertical weight and buoyancy are outside this planar balance. No direction answer is present
\end{definition}
\begin{proof}
  This is the declared model definition of Satisfies Barge Newtons Second Law.
\end{proof}
\begin{definition}[Answer Choice]
  \label{def:physics:phyx-mini-0760:phyxminiproblems-problemphyxmini0760-answerchoice}
  \lean{PhyXMiniProblems.ProblemPhyXMini0760.AnswerChoice}
  Labels of the four direction choices supplied with the problem
\end{definition}
\begin{proof}
  This is the declared model definition of Answer Choice.
\end{proof}
\begin{definition}[displayed Direction Degrees]
  \label{def:physics:phyx-mini-0760:phyxminiproblems-problemphyxmini0760-displayeddirectiondegrees}
  \lean{PhyXMiniProblems.ProblemPhyXMini0760.displayedDirectionDegrees}
  \uses{def:physics:phyx-mini-0760:phyxminiproblems-problemphyxmini0760-answerchoice}
  Degree readout printed beside each displayed answer label
\end{definition}
\begin{proof}
  This is the declared model definition of displayed Direction Degrees.
\end{proof}
\begin{definition}[recorded Dataset Answer]
  \label{def:physics:phyx-mini-0760:phyxminiproblems-problemphyxmini0760-recordeddatasetanswer}
  \lean{PhyXMiniProblems.ProblemPhyXMini0760.recordedDatasetAnswer}
  \uses{def:physics:phyx-mini-0760:phyxminiproblems-problemphyxmini0760-answerchoice}
  The dataset's recorded answer label
\end{definition}
\begin{proof}
  This is the declared model definition of recorded Dataset Answer.
\end{proof}
\begin{definition}[Is Unique Closest Displayed Direction]
  \label{def:physics:phyx-mini-0760:phyxminiproblems-problemphyxmini0760-isuniqueclosestdisplayeddirection}
  \lean{PhyXMiniProblems.ProblemPhyXMini0760.IsUniqueClosestDisplayedDirection}
  \uses{def:physics:phyx-mini-0760:phyxminiproblems-problemphyxmini0760-anglefromdegrees, def:physics:phyx-mini-0760:phyxminiproblems-problemphyxmini0760-answerchoice, def:physics:phyx-mini-0760:phyxminiproblems-problemphyxmini0760-displayeddirectiondegrees}
  A choice is strictly closer on the angle circle than every alternative
\end{definition}
\begin{proof}
  This is the declared model definition of Is Unique Closest Displayed Direction.
\end{proof}
\begin{lemma}[water Force Components]
  \label{lem:physics:phyx-mini-0760:phyxminiproblems-problemphyxmini0760-waterforcecomponents}
  \lean{PhyXMiniProblems.ProblemPhyXMini0760.waterForceComponents}
  \uses{def:physics:phyx-mini-0760:phyxminiproblems-problemphyxmini0760-xaxis, def:physics:phyx-mini-0760:phyxminiproblems-problemphyxmini0760-yaxis, def:physics:phyx-mini-0760:phyxminiproblems-problemphyxmini0760-forcevectorinnewtons, def:physics:phyx-mini-0760:phyxminiproblems-problemphyxmini0760-anglefromdegrees, def:physics:phyx-mini-0760:phyxminiproblems-problemphyxmini0760-canalbargesetup, def:physics:phyx-mini-0760:phyxminiproblems-problemphyxmini0760-matchesproblemreadouts, def:physics:phyx-mini-0760:phyxminiproblems-problemphyxmini0760-satisfiesstraightlinekinematics, def:physics:phyx-mini-0760:phyxminiproblems-problemphyxmini0760-satisfiesplanarforcedirectionlaw, def:physics:phyx-mini-0760:phyxminiproblems-problemphyxmini0760-satisfiesbargenewtonssecondlaw}
  Resolving Newton's second law into the figure's two axes gives the exact SI components of the unknown water force. This is a derived lemma, not a premise field; in particular its negative y component is not figure input
\end{lemma}
\begin{proof}
  The conclusion follows from the governing relations and figure readouts named in the statement of water Force Components.
\end{proof}
% --- Archon declaration topology end ---
% --- Archon physics formalization source end ---
